\documentclass[11pt]{article}
\usepackage[margin=0.9in]{geometry}
\usepackage{amsmath,amssymb,mathtools,microtype}
\usepackage{xcolor}
\usepackage[colorlinks=true,linkcolor=blue!55!black,citecolor=blue!55!black,urlcolor=blue!55!black]{hyperref}

\newcommand{\Las}{\Pi_{\mathrm{as}}}
\newcommand{\Lmult}{\Pi_{\mathrm{mult}}}

\title{Three Open Problems on Summing Multilinear Operators\\Resolved by Later Sharp Theorems}
\author{Literature-implied answer for arXiv:1204.5411\\Model: GPT5.6}
\date{19 July 2026}

\begin{document}
\maketitle

\begin{center}
\fbox{\parbox{0.92\textwidth}{\textbf{Status: literature-implied answer.}
All three problems in Section~4 of Pellegrino's source paper are answered by
combining its transfer formulas with the later optimality theorems
\cite{PS,AP}.  The deductions below are not presented as new theorems.}}
\end{center}

\section{Source questions}

The source paper \cite[p.~8, Section~4]{Source} asks:
\begin{enumerate}
\item whether, for some $\theta\in[0,1]$, the coincidence
\begin{equation}\label{eq:interp}
 \mathcal L({}^n\ell_1;\ell_2)
 =\Las\!\left(\frac{2}{1+\theta};
       \frac{2n}{n+\theta},\ldots,\frac{2n}{n+\theta}\right)
       ({}^n\ell_1;\ell_2)
\end{equation}
is sharp;
\item for real $s>2$ and integer $n\ge s$, the exact value of
\begin{equation}\label{eq:M}
\begin{aligned}
 M_{s,n}:=\sup\bigl\{r:\;&\mathcal L(E_1,\ldots,E_n;\mathbb K)\\
 &=\Lmult(2;r,\ldots,r)(E_1,\ldots,E_n;\mathbb K)\\[-2pt]
 &\text{for all cotype-$s$ spaces }E_j\bigr\};
\end{aligned}
\end{equation}
\item whether $M_{s,n}$ can equal the analogous absolutely-summing
supremum outside the case $(s,n)=(2,2)$.
\end{enumerate}

\section{Problem 1: affirmative, for every \texorpdfstring{$0<\theta\le1$}{0<theta<=1}}

Put $p_\theta=2/(1+\theta)$ and $q_\theta=2n/(n+\theta)$.
The cotype-$2$ transfer identity used in the source gives, for $1\le t\le2$,
\begin{equation}\label{eq:transfer}
 \Las(p_\theta;t,\ldots,t)({}^n\ell_1;\ell_2)
 =\Las\!\left(\rho_\theta(t);1,\ldots,1\right)
 ({}^n\ell_1;\ell_2),
 \qquad
 \rho_\theta(t)=
 \frac{p_\theta t}{n p_\theta(t-1)+t}.
\end{equation}
Here $\rho_\theta(q_\theta)=2/(n+1)$ and $\rho_\theta$ is strictly
decreasing.  Pellegrino and Seoane-Sep\'ulveda proved that $2/(n+1)$ is the
optimal output exponent in
\(\mathcal L({}^n\ell_1;\ell_2)=\Las(r;1,\ldots,1)\)
\cite[Theorem~1.1]{PS}.

If the common input exponent in \eqref{eq:interp} could be increased from
$q_\theta$ to some $t>q_\theta$, take $t$ sufficiently close to $q_\theta$
that $t\le2$.  Equation \eqref{eq:transfer} would give coincidence at
$\rho_\theta(t)<2/(n+1)$, contradicting that theorem.  The same calculation,
with the output exponent decreased slightly from $p_\theta$, again produces
an exponent below $2/(n+1)$.  Thus \eqref{eq:interp} is sharp in both natural
exponent directions for every $0<\theta\le1$.  In particular, Problem~1 has
an affirmative answer.  No claim about the separate endpoint $\theta=0$ is
needed for the source question.

\section{Problems 2 and 3: exact values}

Ara\'ujo and Pellegrino prove that every scalar $n$-linear form on arbitrary
Banach spaces is multiple
\((2;2n/(2n-1),\ldots,2n/(2n-1))\)-summing
\cite[Theorem~3.2(i)]{AP}.  Hence
\[
 M_{s,n}\ge \frac{2n}{2n-1}.
\]
For the reverse inequality, specialize their Corollary~3.1 to multilinearity
$n$, output exponent $2$, and domain exponent $p=s$.  Its hypothesis is
$2\le s<2n$, which follows from $2<s\le n$, and it gives
\[
 \sup\{r:\mathcal L({}^n\ell_s;\mathbb K)
 =\Lmult(2;r,\ldots,r)({}^n\ell_s;\mathbb K)\}
 \le \frac{2n}{2n-1}.
\]
Since $\ell_s$ has cotype $s$, it is an admissible test space in
\eqref{eq:M}.  Therefore
\begin{equation}\label{eq:exactM}
 \boxed{M_{s,n}=\frac{2n}{2n-1}},
\end{equation}
and the supremum is attained.  This answers Problem~2.

The source itself proves that the corresponding absolutely-summing supremum
is attained and equals
\[
 A_{s,n}=\frac{2sn}{2sn+s-2n}
 \qquad\text{\cite[Eq.~(2.10)]{Source}.}
\]
Combining this with \eqref{eq:exactM}, direct cross-multiplication gives
\[
 M_{s,n}=A_{s,n}
 \quad\Longleftrightarrow\quad n=s.
\]
Thus Problem~3 has a complete classification: equality holds for every
integer pair $(s,n)=(k,k)$ with $k>2$, and fails whenever $n>s$.

\section{Provenance and scope}

The supporting paper \cite{AP} explicitly cites page~194 of the published
source and improves the interval containing $M_{s,n}$, but it does not state
the $\ell_s$ specialization above or the equality classification.  The paper
\cite{PS} proves the endpoint optimality used for Problem~1 while identifying
a different named conjecture.  Exact searches did not locate a paper stating
the three deductions together.  Accordingly this packet records a
\emph{literature-implied answer}, not a new proof result.

\begin{thebibliography}{9}
\bibitem{Source}
D. Pellegrino,
\emph{Sharp coincidences for absolutely summing multilinear operators},
Linear Algebra Appl. 440 (2014), 188--196; arXiv:1204.5411.

\bibitem{PS}
D. Pellegrino and J. B. Seoane-Sep\'ulveda,
\emph{Grothendieck's theorem for absolutely summing multilinear operators is optimal},
Linear Multilinear Algebra 63 (2015), 554--558; arXiv:1307.4809.

\bibitem{AP}
G. Ara\'ujo and D. Pellegrino,
\emph{Optimal estimates for summing multilinear operators},
Linear Multilinear Algebra 65 (2017), 930--942; arXiv:1403.6064.
\end{thebibliography}

\end{document}
